%! Piecewise-Defined Functions — Unit 2, Lesson 2.6 — Exit Ticket
\documentclass[10pt]{article}
\usepackage{atda-article}
\usepackage{atda-boxes}
\newcommand{\TallMath}[1]{$\displaystyle #1\rule[-1.4em]{0pt}{3.2em}$}
\setlength{\workrowsep}{8pt}

\begin{document}

\pageheader{Unit 2, Lesson 2.6}{Exit Ticket}
% NAMESTRIP: no \namedateperiod here. The student writes their name once, on the
% cover this component is stapled behind. See references/conventions.md.

\vspace{0.15in}

\begin{notesbox}{One Babysitting Rate, and a Corner Called a Break}
\small
Notes closed. Three items. A babysitter charges a flat \$40 for a job of up to and including $4$
hours, and \$15 for each hour after the fourth. The graph shows $B(h)$, the charge in dollars for a
job of $h$ hours. No job runs longer than $8$ hours.

\begin{center}
\begin{tikzpicture}
\begin{axis}[
    width=10.2cm, height=4.4cm,
    xlabel={$h$ (hours of the job)}, ylabel={$B$ (charge, \$)},
    xmin=0, xmax=8, ymin=0, ymax=110,
    xtick={0,1,2,3,4,5,6,7,8}, ytick={0,20,40,60,80,100}, minor y tick num=1,
    grid=both, grid style={linegray!45}, minor grid style={linegray!30},
    axis lines=left,
    tick label style={font=\tiny}, label style={font=\scriptsize}]
  \addplot[royal, very thick] coordinates {(0,40) (4,40) (8,100)};
\end{axis}
\end{tikzpicture}
\end{center}

\begin{enumerate}[leftmargin=1.6em, itemsep=7pt, topsep=4pt]

\item Give $B(3)$ and $B(6)$, and for each one name which piece owns that input --- the flat piece or
the hourly piece. Then give the \emph{boundary point} of this function.

\writelines{3}

\item Does $B$ have an absolute maximum? Does it have an absolute minimum? For each, give a value
\emph{and} a location. Read the location of the minimum carefully.

\writelines{2}

\item A classmate writes: \emph{``The graph of $B$ has a break at $h=4$, because that is where the
rule changes.''} Naming $h=4$ was correct in item 1. Explain what is wrong with the sentence anyway,
and write the correct statement.

\writelines{3}

\end{enumerate}
\end{notesbox}

\end{document}
